\chapter{الإطار العام للمشروع}

\begin{chaptersummary}
نبين في هذا الفصل تعريفاً بالمشروع والهدف منه إضافة إلى ذكر المتطلبات الوظيفية وغير الوظيفية.
\end{chaptersummary}

\section{مقدمة عامة عن تطور التعليم (\en{Evolution of Education})}
% Content for Chapter 1, Section 1.1
شهدت العملية التعليمية عبر العصور تحولات بنيوية استجابةً للمتغيرات التقنية والاجتماعية؛ فبعد أن كان التعليم يعتمد على التلقين المباشر والمناهج الموحدة التي غالباً ما تُغفل الفروق الفردية، أتاحت التكنولوجيا الرقمية آفاقاً جديدة للتعلم المرن. ومع مطلع القرن الحادي والعشرين، برزت أنظمة إدارة التعلم (\en{LMS}) ومنصات مؤتمرات الفيديو، والتي تسارع تبنيها بشكل كبير إبان جائحة \en{COVID-19}.

إلا أن هذه المنصات الحالية بقيت مجرد ناقل غير تفاعلي للصوت والصورة، حيث يمر المحتوى التعليمي عبرها كتيار بيانات عابر لا يترك أثراً منظماً بعد انتهاء الجلسة، مما يضع عبئاً إدراكياً كبيراً على المعلم في متابعة الفهم، ويترك الطالب في حالة من التلقي غير التفاعلي.

يأتي مشروع \en{IntelliLect} ليمثل الجيل التالي من هذه المنصات، عبر دمج البث المباشر بقدرات الذكاء الاصطناعي التوليديي في الوقت الفعلي. لم يعد الذّكاء الصنعيّ هنا مجرد أداة خارجية، بل "مساعد تدريسي رقمي" يعمل في خلفية المحاضرة؛ فيحول الكلام المنطوق لحظياً إلى نصوص وملخصات، ويضمن دقة المعلومات عبر تقنيات التوليد المعزز بالاسترجاع (\en{RAG})، مما يفرغ المعلم لجودة الشرح ويمنح الطالب تجربة تعليمية تفاعلية ثرية. % This line pulls the text from your new file

\section{الهدف من المشروع}

يهدف المشروع إلى تطوير منصة \en{Intelligent Lecture}) \en{IntelliLect}) وهي منصة تعليمية مدعومة بالذكاء الصنعي للتعلم المتزامن تجمع بين استقرار البث المباشر وقدرات الذكاء الصنعي التوليدي، وذلك لتحويل عملية التدريس عبر الإنترنت من مجرد نقل سلبي للصوت والصورة إلى تجربة تعليمية تفاعلية غنية بالبيانات، تساهم في رفع كفاءة الفهم لدى الطلاب وتخفيف العبء التدريسي عن المعلم.

\section{المتطلبات}

\subsection{المتطلبات الوظيفية (\en{Functional Requirements})}

\noindent المستخدمون غير المسجلين: يجب على النظام أن يسمح للمستخدمين غير المسجلين ب:

\begin{enumerate}
  \item طلب التسجيل في النظام إما كطالب أو كمعلم.
  \item تسجيل الدخول إلى النظام.
  \item إعادة تعيين كلمة المرور.
  \item طلب رمز إعادة تعيين كلمة المرور.
\end{enumerate}

\noindent المستخدمون المصرح لهم: يجب على النظام أن يسمح للمستخدمين المصرح لهم ب:

\begin{enumerate}
  \item رؤية معلومات ملفهم الشخصي.
  \item تحديث معلومات ملفهم الشخصي.
  \item تغيير كلمة المرور.
  \item تسجيل الخروج من النظام.
\end{enumerate}

\noindent مدير النظام الأعلى (\en{Super System Admin}): يجب على النظام أن يسمح للمدير الأعلى ب:

\begin{enumerate}
  \item تسجيل الدخول باستخدام المصادقة الثنائية.
  \item إدارة مدراء النظام:
  \begin{enumerate}
    \item استعراض مدراء النظام الأدنى منه.
    \item إضافة مدير نظام أدنى منه، أو تعطيل حسابه وإعادة تنشيطه.
  \end{enumerate}
  \item إدارة المستخدمين على مستوى المنصة:
  \begin{enumerate}
    \item استعراض جميع المستخدمين والبحث فيهم فلترتهم حسب الدور والحالة وتاريخ الإنشاء.
    \item عرض تفاصيل أي مستخدم وعضويته في الفصول الدراسية.
    \item قبول أو رفض طلبات التسجيل، وتعطيل الحسابات أو إعادة تنشيطها.
  \end{enumerate}
  \item الإشراف على الفصول الدراسية:
  \begin{enumerate}
    \item استعراض جميع الفصول الدراسية والبحث فيها.
    \item إنشاء فصل دراسي أو تعديله أو حذفه.
    \item إسناد معلم الفصل أو تغييره (نقل الملكية).
    \item إضافة طلاب إلى الفصل أو إزالتهم.
    \item معاينة أثر الحذف (عدد الجلسات والملفات والتسجيلات والملخصات المتأثرة) قبل تنفيذه.
  \end{enumerate}
  \item الإشراف على جلسات التعلم:
  \begin{enumerate}
    \item استعراض جميع الجلسات والبحث فيها.
    \item عرض الجلسات النشطة في الوقت الفعلي (عدد المشاركين، حالة التسجيل، حالة المساعد الذكي).
    \item الإنهاء القسري لجلسة مباشرة متوقفة، على أن يتم ذلك بصورة آمنة تضمن اكتمال تسجيل المحاضرة وتوليد ملخصها.
    \item حذف جلسة مع ما يتبعها من تسجيلات ومخرجات.
  \end{enumerate}
  \item إدارة المحتوى وقاعدة المعرفة:
  \begin{enumerate}
    \item استعراض جميع الملفات المرفوعة عبر المنصة وحالة فهرستها.
    \item عرض تفاصيل أخطاء الفهرسة لتشخيصها.
    \item إعادة فهرسة ملف مفرد أو ملفات فصل دراسي كامل.
    \item حذف ملف مع مقاطعه المفهرسة.
    \item استعراض إحصائيات قاعدة المعرفة لكل فصل دراسي وللمنصة ككل.
  \end{enumerate}
  \item إدارة التسجيلات والملخصات:
  \begin{enumerate}
    \item استعراض جميع تسجيلات المحاضرات وملخصاتها وحالاتها.
    \item حذف تسجيل أو ملخص.
  \end{enumerate}
\end{enumerate}

\noindent مدير النظام (\en{System Admin}): يجب على النظام أن يسمح لمدير النظام ب:

\begin{enumerate}
  \item استعراض المستخدمين وطلبات التسجيل المعلقة.
  \item مراجعة طلبات التسجيل المقدمة من المعلمين والطلاب وقبولها أو رفضها.
  \item إلغاء تنشيط حسابات الطلاب/المعلمين الحالية وإعادة تنشيطها.
\end{enumerate}

\noindent المعلم: يجب على النظام أن يسمح للمعلم ب:

\begin{enumerate}
  \item إدارة فصوله الدراسية:
  \begin{enumerate}
    \item عرض الفصول الدراسية المسجل فيها
    \item إرفاق ملفات بالفصل الدراسي.
    \item تعديل معلومات فصوله الدراسية.
  \end{enumerate}
  \item إدارة جلسات التعلم:
  \begin{enumerate}
    \item بدء جلسة تعلم في الوقت الفعلي (بث مباشر).
    \item إرسال ملاحظات في دردشة الجلسة في الوقت الفعلي.
    \item رؤية طلابه أثناء التدريس.
    \item استخدام السبورة البيضاء \en{white board} للشرح.
    \item معرفة عدد المشاركين في الجلسة، والتحكم بالسماح للطلاب بمشاركة الصوت أو الصورة.
    \item تحسين أسلوبه التدريسي من خلال تلقي إشعارات خاصة من الذكاء الصنعي في الوقت الفعلي.
  \end{enumerate}
  \item إدارة الاختبارات (\en{Quizzes}) داخل الجلسة:
  \begin{enumerate}
    \item إنشاء مسودة اختبار يدويًّا أو توليدها من نص الجلسة، وتعديلها قبل نشرها.
    \item توليد سؤال إضافي، أو توليد خيارات الإجابة لسؤال كتبه بنفسه.
    \item نشر الاختبار وإغلاقه أو إلغاؤه.
    \item تمديد مهلة الإجابة للصف كاملًا أو لطلاب بعينهم.
    \item استعراض نتائج الاختبار وإجابات الطلاب.
    \item متابعة تحصيل الطلاب في الفصل الدراسي عبر جلساته كلّها.
  \end{enumerate}
\end{enumerate}

\noindent الطالب: يجب على النظام أن يسمح للطالب ب:

\begin{enumerate}
  \item التفاعل مع جلسات الفصل الدراسي:
  \begin{enumerate}
    \item الانضمام إلى جلسة تعلم.
    \item رؤية شرح المعلم على السبورة البيضاء.
    \item طرح الأسئلة نصيًّا أثناء الجلسة، والتحدّث صوتيًّا عندما يسمح المعلم بمشاركة الصوت.
  \end{enumerate}
  \item التعامل مع الاختبارات:
  \begin{enumerate}
    \item الإجابة عن أسئلة الاختبار المفتوح ضمن المهلة المحدّدة، وتسليمه.
    \item الاطّلاع على نتيجته بعد إغلاق الاختبار.
    \item متابعة تحصيله في الفصل الدراسي ومقارنته بمتوسّط الصف.
  \end{enumerate}
\end{enumerate}

\noindent النظام: يجب أن يكون النظام قادرا على:

\begin{enumerate}
  \item إدارة جلسة تدريس في الوقت الفعلي بين المعلم وعدد من الطلاب.
  \item إرسال إشعارات في الوقت الفعلي لجميع الأطراف (الطلاب والمعلمين).
  \item إجراء تلخيص، وتوليد كلمات مفتاحية، وتحديد النقط التي تم التركيز عليها في المحاضرة باستخدام الملفات التي رفعها المعلم والنص المستخرج من صوت المعلم أثناء التدريس.
  \item حفظ تسجيل المحاضرة.
  \item توليد اختبارات سياقية أثناء الجلسة بناء على النص المباشر للمحاضرة باستخدام النماذج اللغوية الكبيرة (\en{LLM})، على أن يراجعها المعلم قبل نشرها على الطلاب.
  \item تحليل النص المباشر للمعلم مقابل الملفات المرجعية المحملة (ملفات \en{PDF} مثلا)باستخدام تقنية التوليد المعزز بالاسترجاع (\en{RAG})
  \item يجب أن يكتشف النظام التناقضات أو التفسيرات غير الواضحة ويرسل اقتراحا خاصا في الوقت الفعلي للمعلم.
  \item حفظ جلسة التدريس حتى تتمكن الأطراف من تحميلها لمشاهدتها لاحقا.
  \item يجب أن يبني النظام قاعدة معرفية قابلة للبحث من الملفات التي يرفعها المعلم، وذلك عبر استخراج نصوصها وتقسيمها إلى مقاطع وفهرستها.
  \item يجب أن يسترجع النظام المقاطع الأكثر صلة باستعلام معيّن، على أن يكون البحث مقصورًا على مواد الفصل الدراسي المعني دون سواه.
  \item تصحيح إجابات الاختبارات آليًّا على الخادم واحتساب العلامات.
  \item منع وصول مفتاح الإجابة إلى الطالب، وحجب نتيجته حتى إغلاق الاختبار.
  \item إعلام المشاركين بتغيّر حالة الاختبار في الوقت الفعلي.
  \item استبعاد الاختبار الملغى من احتساب العلامات مع الاحتفاظ بإجابات الطلاب.
\end{enumerate}

\subsection{المتطلبات غير الوظيفية (\en{Non-Functional Requirements})}

\begin{enumerate}
  \item يجب ألا يتجاوز زمن استجابة (\en{Latency}) بث الفيديو/الصوت 500 مللي ثانية.
  \item يجب أن تكون بنية النظام قابلة للتوسع أفقيا، مع قدرة النموذج الأولي الحالي على التعامل مع إلى جلسات متزامنة (تصل إلى 50 مستخدما نشطا) على أجهزة التطوير المحلية دون توقف أو انهيار النظام.
  \item يجب ألا يتجاوز زمن استجابة المساعد الذكي 3 إلى 5 ثوا ن بعد انتهاء المعلم من شرح مفهوم معين.
\end{enumerate}
